\chapter{Personen}

\begin{itemize}
\item Aristoteles: Zeit als Zahl der Bewegung; Richtungen nur in beseelten Körpern; Ort als Gefäß. Kapitel 4, 5, 10.
\item Baumann, Julius: Kritik der Parallelisierung; drei Zeiten; Ich-Dauer; Zahl aus dem Geist. Kapitel 0, 2, 4, 7, 8, 10.
\item Bonsiepen, Wolfgang: Hegels Raum-Zeit-Lehre und ihre Genese; Edition der Nachschrift. Kapitel 1, 2, 5, 6, 8, 9.
\item Brentano, Franz: Kontinuum als Ganzes und als Grenze; sachliche und modale Differenzen; Überordnung des Zeitlichen. Kapitel 2, 5, 6, 7, 10.
\item Chiba, Kiyoshi: raumzeitliche Gegenstände; Warten und Gehen; zeit-transzendenter Standpunkt. Kapitel 6, 11.
\item Einstein, Albert: Relativitätstheorie, nach Schneider und Reichenbach. Kapitel 1, 5.
\item Fink, Eugen: das Und; Zenons Parallelisierung; Jetzt weltweit, Hier je meines; Räumen und Zeitigen. Kapitel 0, 1, 2, 3, 5, 9, 11.
\item Guzzoni, Ute: Weile und Weite; Ort als Augenblick des Raumes; die Zeit wird zum Raum. Kapitel 0, 2, 5, 8, 9, 10.
\item Hegel, Georg Wilhelm Friedrich: zwei Weisen des Außersichseins; die Zeit als wahrhafter Punkt; Ort als räumliches Jetzt; Zeit und Ich = Ich. Kapitel 0, 1, 2, 5, 6, 7, 8, 9.
\item Heidegger, Martin: Zeit-Raum; Rücknahme der Rückführung der Räumlichkeit auf die Zeitlichkeit, nach Guzzoni und Fink. Kapitel 0, 5, 9, 11.
\item Hölderlin, Friedrich: \zitat{Zeit eilt hin zum Ort}, nach Guzzoni. Kapitel 5.
\item Kant, Immanuel: Anschauungsformen; Schematismus; Widerlegung des Idealismus; inkongruente Gegenstücke; Dimensionszahl aus dem Kraftgesetz. Kapitel 1, 4, 7, 8, 11.
\item Koch, Anton Friedrich: Spezifizierbarkeit und Indexikalität; Richtungen aus dem Leib; Physik als Abstraktion von der Subjektivität. Kapitel 2, 3, 4, 6, 7, 9.
\item Leibniz, Gottfried Wilhelm: Raum als Ordnung der Koexistenzen, nach Baumann. Kapitel 8.
\item Martić, Marko: Kants synthetischer Vorrang der Zeit; Zeitraum; Newtons Anisotropie der Zeit. Kapitel 1, 4, 7, 8.
\item Minkowski, Hermann: Zeit als imaginäre Koordinate, nach Schneider und Reichenbach. Kapitel 1, 10.
\item Newton, Isaac: Anisotropie der Zeit, Isotropie des Raumes, nach Martić. Kapitel 4.
\item Platon: Chora und Idee, nach Fink. Kapitel 9.
\item Reichenbach, Hans: Zuordnungsdefinition; Kausaltheorie der Zeit; Faserrichtung; Lichtgeometrie; Dimensionszahl; ewiges Jetzt. Kapitel 1, 3, 4, 5, 6, 7, 8, 9, 10, 11.
\item Schelling, Friedrich Wilhelm Joseph: wechselseitige Endlichkeit von Zeit und Raum, nach Martić. Kapitel 8.
\item Schneider, Ilse: Raum und Zeit verlieren die physikalische Gegenständlichkeit; Koordinaten als Symbole. Kapitel 1, 4, 10.
\item Wagner, Richard: \zitat{Zum Raum wird hier die Zeit}, nach Guzzoni. Kapitel 5.
\item Zenon von Elea: Paradoxien als Parallelisierung, nach Fink. Kapitel 1, 8.
\end{itemize}
